\documentclass[11pt]{article}

\usepackage[a4paper,margin=1in]{geometry}
\usepackage{amsmath,amssymb,amsthm}
\usepackage{booktabs}
\usepackage{hyperref}
\usepackage{xcolor}
\usepackage{microtype}

\newcommand{\Vsix}{V_{600}}
\newcommand{\Vtwentyfour}{V_{24}}
\newcommand{\Lsix}{L_{600}}
\newcommand{\Esix}{E_{600}}
\newcommand{\Elift}{E_{\mathrm{lifted}}}
\newcommand{\Eres}{E_{\mathrm{residual}}}
\newcommand{\twoI}{2I}
\newcommand{\twoT}{2T}

\theoremstyle{definition}
\newtheorem{definition}{Definition}

% Avoid TeX-inserted hyphenated line breaks rendering as artifacts in some
% viewers; keep breaks at existing hyphens.
\hyphenpenalty=10000
\exhyphenpenalty=100
\setlength{\emergencystretch}{3em}

\title{From Schl\"afli Decomposition to Spectral Bridge:\\
\large A First Explicit Closure-Projection Channel in $\Vsix$}
\author{%
  Lee Smart\\[4pt]
  \textit{Institute of Vibrational Field Dynamics}\\[2pt]
  \texttt{contact@vibrationalfielddynamics.org}\\[2pt]
  \texttt{@vfd\_org}%
}
\date{May 2026}

\begin{document}
\maketitle

\begin{center}
\fbox{\parbox{0.92\textwidth}{\small
\textbf{Status:} Synthesis note / programme map. Not peer-reviewed.

\medskip
This note introduces no new load-bearing mathematical claim.  It
synthesises two previously published reproducible exact-arithmetic
notes --- the Schl\"afli decomposition of the 600-cell~\cite{SmartSchlafli2026}
and the 24--600 spectral bridge~\cite{SmartBridge2026} --- and
explains their combined role as a first local-to-global spectral
channel in the $\Vsix$ programme.  Both prior notes are
self-contained and reproducible from public repositories; this note
is a reader's map between them.
}}
\end{center}

\bigskip

\begin{abstract}
The Schl\"afli decomposition of the 600-cell represents $\Vsix$ as
the binary icosahedral group $\twoI$ and partitions it into five
right cosets of the binary tetrahedral subgroup $\twoT$, each
carrying intrinsic 24-cell structure.  The 24--600 spectral bridge
then shows that the local $\lambda=12$ eigenspaces of these five
24-cell sectors zero-extend exactly into the global $\lambda=12$
eigenspace of the 600-cell Laplacian.  This note synthesises those
two exact-arithmetic results and explains their combined role as a
first explicit local-to-global spectral channel in the $\Vsix$ programme.
No new physical claim is made.
\end{abstract}

\section{Why this synthesis note exists}\label{sec:why}

Two reproducible exact-arithmetic notes have been published.  Read
individually, each is narrow:

\begin{itemize}
  \item the Schl\"afli note~\cite{SmartSchlafli2026} proves the
        five-sector finite geometry of the 600-cell;
  \item the 24--600 spectral bridge~\cite{SmartBridge2026} proves
        a single $\lambda=12$ local-to-global spectral lift.
\end{itemize}

Read together, they exhibit a pattern: a local shell structure
inside the 600-cell does not merely sit inside the global geometry
--- at a specific spectral value it couples \emph{exactly} to a
global invariant sector.  The pattern is short and precise enough to
warrant a brief synthesis.

This note is a map, not a new claim-heavy paper.  Section~\ref{sec:dec}
summarises the Schl\"afli decomposition; Section~\ref{sec:bridge}
summarises the spectral bridge; Section~\ref{sec:nontrivial} explains
why their combination is non-trivial; Section~\ref{sec:channel}
defines the term ``closure-projection channel'' modestly;
Section~\ref{sec:scope} states what is and is not established;
Section~\ref{sec:programme} explains the role of the synthesis in the
broader $\Vsix$ programme; Section~\ref{sec:repro} lists the
dependency map and reproduction status.

\section{Artifact I: the five 24-cell cosets}\label{sec:dec}

The Schl\"afli note~\cite{SmartSchlafli2026} establishes, by exact
$\mathbb{Q}(\sqrt 5)$ arithmetic, that:

\[
  \Vsix \;\cong\; \twoI, \qquad \twoT \subset \twoI, \qquad
  [\twoI : \twoT] = 5,
\]
\[
  \Vsix \;=\; C_0 \sqcup C_1 \sqcup C_2 \sqcup C_3 \sqcup C_4,
  \quad C_i = \twoT\cdot g_i, \quad |C_i| = 24.
\]
Each coset $C_i$ is an isometric copy of the 24-cell as a metric
subspace of $S^3$.  Two further structural facts are used below:

\begin{itemize}
  \item \textbf{Two distinct distance shells.}  Inside $\Vsix$, the
        shortest non-zero squared distance is
        $d^2_{600} = (3 - \sqrt 5)/2 \approx 0.382$ (the 600-cell
        edge); inside each coset $C_i$, the shortest non-zero
        squared distance between members is the next shell up,
        $d^2_{24} = 1$.  The two graphs share \emph{no} edges: no
        intra-coset pair is at separation $d^2_{600}$, and every
        600-cell edge runs between two different cosets.
  \item \textbf{Coset action.}  Right multiplication by $\twoI$
        permutes the five cosets.  The kernel of this action is
        exactly $\{\pm 1\}\subset\twoT$, and the image is $A_5 =
        \twoI/\{\pm 1\}$ acting transitively on five points.
\end{itemize}

All three facts are certified by exact rational arithmetic; see
\cite{SmartSchlafli2026} for the proofs and the reproduction
instructions.

\section{Artifact II: the $\lambda=12$ spectral bridge}\label{sec:bridge}

For each coset $C_i$ let $L_{C_i}$ be the local 24-cell shell
Laplacian (the unweighted combinatorial Laplacian of the
$d^2_{24}=1$ graph on $C_i$), and let $\Lsix$ be the global
600-cell shell Laplacian (the unweighted combinatorial Laplacian of
the $d^2_{600}$ graph on $\Vsix$).  The bridge
note~\cite{SmartBridge2026} establishes, by exact rational
arithmetic, the following \emph{lift identity}:

\begin{equation*}
  \mathrm{lift}_i\bigl(E_{C_i}(12)\bigr) \;\subset\; \Esix(12),
\end{equation*}
where $\mathrm{lift}_i\colon\mathbb{R}^{C_i}\to\mathbb{R}^{\Vsix}$
is zero-extension and $E_X(\lambda)$ denotes the
$\lambda$-eigenspace of $L_X$.

Stacking across the five cosets,
\[
  \Elift \;:=\; \bigoplus_{i=0}^{4}
  \mathrm{lift}_i\bigl(E_{C_i}(12)\bigr),
\]
\[
  \dim\Elift = 10, \qquad \dim\Esix(12) = 25.
\]
Under the $A_5$ action of Section~\ref{sec:dec} (after verifying
that $\{\pm 1\}$ acts trivially on $\Esix(12)$),
\[
  \Elift \cong 2\!\cdot\!Y_5, \qquad
  \Eres \cong 3\!\cdot\!Y_5, \qquad
  \Esix(12) \cong 5\!\cdot\!Y_5,
\]
with $Y_5$ the 5-dimensional irreducible representation of $A_5$.

The bridge is also \emph{selective}: $\lambda = 12$ is the only
eigenvalue of $L_{C_i}$ that admits a full zero-extension.  At
$\lambda = 0$ only the global constant (the symmetric sum of the
five coset indicators) survives; at $\lambda \in \{4, 8, 10\}$
these are not eigenvalues of $\Lsix$ at all (certified exactly),
so the question is vacuous.

\section{Why the bridge is non-trivial}\label{sec:nontrivial}

Sections~\ref{sec:dec} and~\ref{sec:bridge} use \emph{different
Laplacians on different graphs}.  The local Laplacian $L_{C_i}$ is
built from the 24-cell edge graph on 24 vertices at squared
distance $d^2_{24}=1$.  The global Laplacian $\Lsix$ is built from
the 600-cell edge graph on 120 vertices at squared distance
$d^2_{600}=(3-\sqrt 5)/2$.  These are not subgraphs of each other:
no edge of $L_{C_i}$ is an edge of $\Lsix$, and conversely no
$\Lsix$-edge joins two members of the same coset.

In particular, for a function $v$ supported on a single coset
$C_i$, the value $(\Lsix\, v)(u)$ at a vertex $u\notin C_i$ is in
general non-zero --- it is minus the sum of $v$ over the global
neighbours of $u$ that happen to lie in $C_i$.  The zero-extension
of $v$ is therefore generically \emph{not} a $\Lsix$-eigenvector.

The content of the lift identity is that, on the local
$\lambda=12$ eigenspace specifically, all such off-coset
couplings cancel \emph{identically}: a finite-rank linear
condition that turns out to be satisfied at $\lambda=12$ and
nowhere else in the local spectrum.

\medskip
\noindent
Put compactly: \emph{the bridge is not a restriction theorem; it
is a cancellation theorem.}  Equivalently, the $\lambda=12$ lift
is a finite-dimensional compatibility condition between a local
shell operator and a global shell operator built from different
distance shells.

\section{Closure-projection channel: a careful definition}\label{sec:channel}

\begin{definition}[Closure-projection channel, working]\label{def:cpc}
A \emph{closure-projection channel} (in the sense of the present
$\Vsix$ programme) is a reproducible local-to-global compatibility
relation: a class of modes defined on a local geometric sector
projects, lifts, or embeds exactly into a global invariant sector
of a larger geometry.
\end{definition}

Specialised to the 24--600 case:

\begin{quote}
The local $\lambda=12$ eigenspaces on the five 24-cell sectors of
the 600-cell zero-extend exactly into the global $\lambda=12$
eigenspace of $\Lsix$, and that lift is $A_5$-stable.
\end{quote}

Definition~\ref{def:cpc} is \emph{interpretive}.  The formal
mathematical content is the combination of the two exact results
cited in Sections~\ref{sec:dec} and~\ref{sec:bridge}: the
five-coset partition and the $\lambda=12$ lift identity with its
$A_5$ decomposition.  The phrase ``closure-projection channel''
adds no new theorem; it is a name for the pattern.

\section{What is and is not established}\label{sec:scope}

\noindent\textbf{Established (by the two cited artifacts):}
\begin{itemize}
  \item $\Vsix$ decomposes into five 24-cell cosets (exact).
  \item The induced coset action descends to $A_5$ (exact).
  \item Local 24-cell shell operators and the global 600-cell
        shell operator are distinct objects on different distance
        shells (exact).
  \item The local $\lambda=12$ eigenspaces zero-extend into the
        global $\lambda=12$ sector (exact, $\mathbb{Q}$-rational).
  \item The lifted subspace $\Elift$ is $A_5$-stable and
        decomposes as $2\!\cdot\!Y_5$ (exact characters).
\end{itemize}

\noindent\textbf{Not established (in this note or in the two cited
artifacts):}
\begin{itemize}
  \item No derivation of particle masses, the Standard Model, or
        any specific physical observable.
  \item No derivation of quantum mechanics or General Relativity.
  \item No claim that $\lambda=12$ is physically realised, or that
        the 600-cell is a physical structure.
  \item No claim that all closure projection is explained by this
        example.
  \item No proof of the Riemann hypothesis or any number-theoretic
        statement.
\end{itemize}

The synthesis is mathematical and interpretive only.

\section{Why this matters for the $\Vsix$ programme}\label{sec:programme}

The $\Vsix$ programme requires explicit mechanisms by which local
finite-geometric structures can participate in global invariant
sectors.  The Schl\"afli decomposition supplies local sectors; the
spectral bridge supplies an exact compatibility channel.  Together,
they provide a finite test case for the broader idea of
\emph{closure projection}.

The next mathematical step is not to identify this channel with a
physical observable, but to formalise the general selection
principle: \emph{given a local sector, a local operator, a global
operator, and a projection/lift map, when does a local eigenspace
contribute to a global invariant sector?}  The $\lambda=12$ lift in
the 24--600 case is one positive instance and one negative-control
template (at every other local eigenvalue) for that general
question.

\section{Reproducibility and dependency map}\label{sec:repro}

This synthesis note introduces no new computation, so it carries no
new formal certificate of its own.  The dependency on the two prior
artifacts is direct:

\medskip
\noindent
\begin{tabular}{p{0.30\textwidth}p{0.40\textwidth}p{0.22\textwidth}}
\toprule
Artifact & Role & Repro status \\
\midrule
Schl\"afli decomposition~\cite{SmartSchlafli2026}
  & Proves the five 24-cell cosets and the $A_5$ coset action
  & exact $\mathbb{Q}(\sqrt 5)$ certificate \\
\addlinespace
24--600 spectral bridge~\cite{SmartBridge2026}
  & Proves the $\lambda=12$ local-to-global lift identity
  & exact $\mathbb{Q}$-rational certificate \\
\addlinespace
This synthesis note
  & Explains the combined interpretation
  & no new formal certificate \\
\bottomrule
\end{tabular}
\medskip

Both prior artifacts are public:
\begin{itemize}
  \item \url{https://github.com/vfd-org/the-24-600-spectral-bridge}
        --- run \texttt{pytest closure\_transform\_engine/tests/}
        from a clone to re-assert every certificate-level claim of
        both prior notes.
\end{itemize}

\smallskip
\noindent
Although named after the spectral-bridge artifact, this repository
now serves as the public $\Vsix$ foundation bundle for all three
notes:
\begin{itemize}
  \item Paper 1 (Schl\"afli decomposition): \texttt{docs/}\allowbreak\texttt{01-schlafli-decomposition/}
  \item Paper 2 (spectral bridge): \texttt{docs/}\allowbreak\texttt{02-spectral-bridge/}
  \item Paper 3 (this synthesis): \texttt{docs/}\allowbreak\texttt{03-closure-projection-channel/}
\end{itemize}

\bigskip
\noindent
\textbf{End of synthesis.}  The next paper in this programme will
attempt to formalise Definition~\ref{def:cpc} into a general
selection principle, but is out of scope here.

\begin{thebibliography}{9}
\bibitem[Smart 2026a]{SmartSchlafli2026}
L.~Smart, \emph{The Schl\"afli Decomposition of the 600-cell:
Five 24-cell Cosets and the Induced $A_5$ Action.}
Public reproducible note, Institute of Vibrational Field Dynamics,
2026.  Repository:
\url{https://github.com/vfd-org/the-24-600-spectral-bridge}
(folder \texttt{docs/01-schlafli-decomposition/}).
\bibitem[Smart 2026b]{SmartBridge2026}
L.~Smart, \emph{The 24--600 Spectral Bridge: A selective $\lambda=12$
eigenspace embedding from five 24-cell cosets into the 600-cell.}
Public reproducible note, Institute of Vibrational Field Dynamics,
2026.  Repository:
\url{https://github.com/vfd-org/the-24-600-spectral-bridge}
(folder \texttt{docs/02-spectral-bridge/}).
\bibitem[Coxeter 1973]{Coxeter1973}
H.\,S.\,M.~Coxeter, \emph{Regular Polytopes}, 3rd ed., Dover, 1973.
\bibitem[Conway--Sloane 1999]{ConwaySloane1999}
J.\,H.~Conway and N.\,J.\,A.~Sloane, \emph{Sphere Packings, Lattices
and Groups}, 3rd ed., Springer-Verlag, 1999.
\bibitem[Conway et al.\ 1985]{ConwayAtlas1985}
J.\,H.~Conway, R.\,T.~Curtis, S.\,P.~Norton, R.\,A.~Parker and
R.\,A.~Wilson, \emph{Atlas of Finite Groups}, Oxford University
Press, 1985.
\end{thebibliography}

\end{document}
